\section*{الفصل \ref{ch:b1:predominant-reaction} --- توازنات حمض--قاعدة وطريقة التفاعل الغالب}

\begin{solution}{exo:b1:predominant-reaction:1}
القواعد المرافقة: \ce{HSO4-}، \ce{CO3^{2-}}، \ce{NH3}، \ce{OH-}. والأحماض
المرافقة: \ce{NH4+}، \ce{H2PO4-}، \ce{H2O}، \ce{H3O+}. \omterm{def:b1:predominant-reaction:ampholyte}{والأمفوليتات}: \ce{HCO3-}،
\ce{H2O}، \ce{HPO4^{2-}} (و\ce{HSO4-}، قاعدة حمض الكبريتيك وحمض
مزدوجته الثانية).
\end{solution}

\begin{solution}{exo:b1:predominant-reaction:2}
\babelsublr{\ce{HF} ($K_a = \num{6.9e-4}$) $>$ \ce{CH3COOH} (\num{1.7e-5}) $>$ \ce{HClO}
(\num{2.8e-8}) $>$ \ce{NH4+} (\num{5.6e-10})}.
\end{solution}

\begin{solution}{exo:b1:predominant-reaction:3}
\omterm{def:b1:predominant-reaction:strong-acid}{الحمض القوي}: $\mathrm{pH} = 2.00$. \omterm{def:b1:predominant-reaction:strong-acid}{القاعدة القوية}: $[\ce{OH-}] = 10^{-2}$،
$\mathrm{pH} = 14.00 - 2.00 = 12.00$.
\end{solution}

\begin{solution}{exo:b1:predominant-reaction:4}
\ce{CH3COOH} تحت pH 4.76، و\ce{CH3COO-} فوقه. عند pH 3 يغلب الشكل
الحمضي (98\,\%)؛ وفي الدم (pH 7.4) يغلب أيون الإيثانوات. وعند pH 3.76 يكون
كسر الحمض $1/(1 + 10^{-1}) = 0.91$.
\end{solution}

\begin{solution}{exo:b1:predominant-reaction:5}
\omterm{def:b1:predominant-reaction:predominant-reaction}{التفاعل الغالب} \ce{NH3 + H2O <=> NH4+ + OH-}، $K = 10^{-(14.00 -
9.25)} = 10^{-4.75}$. فإن لم يتفاعل من \ce{NH3} إلا قليل: $[\ce{OH-}] = \sqrt{K C} =
\qty{1.3e-3}{mol/L}$، $\mathrm{pH} = 11.13$. التحقق: $[\ce{NH4+}] = 1.3\,\%$
من $C$، و$\mathrm{pH} > \mathrm{p}K_a + 1$: الفرضية صالحة.
\end{solution}

\begin{solution}{exo:b1:predominant-reaction:6}
$\mathrm{pH} = \frac12(4.76 + 2.00) = 3.38$ (والقيمة الدقيقة 3.39)؛
$\alpha = 10^{-3.38}/0.010 = 0.042$: متفكك بنسبة 4\,\%، أي دون 10\,\%،
و$\mathrm{pH} < \mathrm{p}K_a - 1$: الفرضية صالحة.
\end{solution}

\begin{solution}{exo:b1:predominant-reaction:7}
\ce{HF + OH- -> F- + H2O}، $K = 10^{14.00 - 3.16} = 10^{10.84}$: تام. ويحوي
\omterm{def:b1:predominant-reaction:predominant-reaction}{المحلول المكافئ} \qty{0.050}{mol/L} من كل من \ce{HF} و\ce{F-}، وهو
\omterm{def:b1:predominant-reaction:buffer}{محلول منظم}: $\mathrm{pH} = \mathrm{p}K_a = 3.16$.
\end{solution}

\begin{solution}{exo:b1:predominant-reaction:8}
$\mathrm{pH} = 4.76$. يحوّل حمض الكلوريدريك \qty{0.010}{mol} من
الإيثانوات إلى حمض: $\mathrm{pH} = 4.76 + \log(0.040/0.060) = 4.58$، أي
تغير قدره 0.18. وفي الماء النقي كان pH سيهبط من 7.00 إلى 2.00.
\end{solution}

\begin{solution}{exo:b1:predominant-reaction:9}
في الماء يتفاعل الحمضان كلاهما تفاعلًا تامًا ليعطيا \ce{H3O+}: فالمذيب
يسوّي بينهما. وحمض الإيثانويك قاعدة أضعف بكثير من الماء، فلا يتفاعل
فيه الحمضان إلا جزئيًا، وبدرجتين مختلفتين، فيمكن مقارنة
قوتيهما. وأقوى حمض يمكن أن يوجد في الماء هو \ce{H3O+}.
\end{solution}

\begin{solution}{exo:b1:predominant-reaction:10}
\ce{CO3^{2-} + H2O <=> HCO3- + OH-}، $K = 10^{-(14.00 - 10.33)} = 10^{-3.67}$؛
$[\ce{OH-}] = \sqrt{10^{-3.67} \times 0.10} = \qty{4.6e-3}{mol/L}$ (4.6\,\%
من $C$)، $\mathrm{pH} = 11.67$. أما القاعدية الثانية، \ce{HCO3- + H2O <=>
CO2 + H2O + OH-}، فثابتها $K = 10^{-7.63}$: وهي تعطي $[\ce{CO2}] \approx
10^{-7.6}$، أي قيمة مهملة.
\end{solution}

\begin{solution}{exo:b1:predominant-reaction:11}
$[\mathrm{A}] = [\ce{H3O+}] = \alpha C$ و$[\mathrm{HA}] = (1 - \alpha)C$،
فيكون $K_a = \alpha^2C/(1 - \alpha)$. ومع $C = 10^{-3}$: $\alpha^2 +
0.0174\,\alpha - 0.0174 = 0$، $\alpha = 0.12$. فالحمض متفكك بنسبة 12\,\%:
وفرضية الحمض القليل التفكك ($\alpha <
10\,\%$) تخفق، ويجب حل المعادلة التربيعية.
\end{solution}

\begin{solution}{exo:b1:predominant-reaction:12}
يحدد حمض الكلوريدريك $[\ce{H3O+}] = \qty{0.010}{mol/L}$: pH 2.00. ومن ثم
\[
  [\ce{CH3COO-}] = \frac{K_a[\ce{CH3COOH}]}{[\ce{H3O+}]} = \frac{1.74 \times 10^{-5}
  \times 0.10}{0.010} = \qty{1.7e-4}{mol/L},
\]
وهو مهمل بجانب 0.010: فأيونات
\ce{H3O+} الآتية من \omterm{def:b1:predominant-reaction:strong-acid}{الحمض القوي} تزيح توازن \omterm{def:b1:predominant-reaction:strong-acid}{الحمض الضعيف}
إلى الخلف.
\end{solution}

\begin{solution}{pb:b1:predominant-reaction:1}
\textbf{1.} $\ce{H3PO4}/\ce{H2PO4-}$، $\ce{H2PO4-}/\ce{HPO4^{2-}}$،
$\ce{HPO4^{2-}}/\ce{PO4^{3-}}$؛ مثلًا \ce{H3PO4 <=> H2PO4- + H+}.
\textbf{2.} الحدود عند 2.15 و7.21 و12.34.
\textbf{3.} pH 1: \ce{H3PO4}؛ 5: \ce{H2PO4-}؛ 9: \ce{HPO4^{2-}}؛ 13:
\ce{PO4^{3-}}.
\textbf{4.} \ce{H2PO4-} و\ce{HPO4^{2-}}.
\textbf{5.} كل بروتون يترك أيونًا أعلى شحنةً سالبة، يمسك
البروتون التالي بقوة أكبر.
\textbf{6.} \ce{H3PO4 + H2O <=> H2PO4- + H3O+}؛ والصيغة
$\frac12(\mathrm{p}K_a + \mathrm{p}C) = 1.58$ تخالف $\mathrm{pH} <
\mathrm{p}K_a - 1 = 1.15$: فنحو خُمس الحمض يتفاعل.
\textbf{7.} $x^2 + K_{a1}x - K_{a1}C = 0$ مع $K_{a1} = 10^{-2.15}$: $x =
\qty{0.0233}{mol/L}$، $\mathrm{pH} = 1.63$.
\textbf{8.} \ce{2H2PO4- <=> H3PO4 + HPO4^{2-}}، $K = 10^{2.15 - 7.21} =
10^{-5.06}$.
\textbf{9.} $\mathrm{pH} = \frac12(2.15 + 7.21) = 4.68$.
\textbf{10.} $\mathrm{pH} = \frac12(7.21 + 12.34) = 9.78$.
\textbf{11.} لا شيء منها: 1.63 و4.68 و9.78؛ ويلزم خليط من اثنين
منها.
\textbf{12.} \ce{H3PO4 + OH- -> H2PO4- + H2O}، $K = 10^{14.00 - 2.15} =
10^{11.85}$: كمّي.
\textbf{13.} \qty{0.100}{mol/L} من \ce{H2PO4-}: pH 4.68.
\textbf{14.} تحوّل الكمية التالية \qty{0.050}{mol} من الهيدروكسيد نصفه:
\qty{0.050}{mol/L} من \ce{H2PO4-} ومن \ce{HPO4^{2-}}، pH 7.21.
\textbf{15.} \qty{0.100}{mol/L} من \ce{HPO4^{2-}}: pH 9.78.
\textbf{16.} يرتفع pH من 1.63، ويقفز حول $n = 0.100$ (إلى 4.68)،
ويرتفع ببطء مارًّا بالقيمة 7.21 عند $n = 0.150$، ويقفز حول $n = 0.200$ (إلى
9.78)، ويتجاوز 12.34 قرب $n = 0.250$ ويبلغ نحو 12.6 عند $n = 0.300$ (\qty{0.100}{mol/L} من \ce{PO4^3-}: $x^2/(0.100 - x) = 10^{-1.66}$، $x = 0.037$، pH 12.57).
\textbf{17.} يحوي كميتين متساويتين من حمض مزدوجة وقاعدتها
ذات $\mathrm{p}K_a$ 7.21 (\cref{prop:b1:predominant-reaction:buffer}).
\textbf{18.} $10^{7.20 - 7.21} = 0.977$.
\textbf{19.} $n(\ce{H2PO4-}) = 0.100 - x$، $n(\ce{HPO4^{2-}}) = x$.
\textbf{20.} $x/(0.100 - x) = 0.977$، $x = \qty{0.0494}{mol}$.
\textbf{21.} يحوّل الحمض \qty{0.0010}{mol} من \ce{HPO4^{2-}}:
$\mathrm{pH} = 7.21 + \log(0.0484/0.0516) = 7.18$، أي تغير قدره $-0.02$. وفي
الماء النقي: من 7.00 إلى 3.00.
\textbf{22.} لا: يتعلق pH بنسبة الكميتين، التي لا يغيّرها
التخفيف (ما دامت النشاطيات قريبة من
التراكيز).
\textbf{23.} $0.100 + 0.0494 = \qty{0.149}{mol}$ من الهيدروكسيد: \textbf{\qty{149}{mL}}
من المحلول ذي التركيز \qty{1.00}{mol/L}، ثم يُكمَل الحجم إلى \qty{1.00}{L}.
\end{solution}
